\chapter{固有态射的基变换定理：证明的完成}
%\addtocontents{toc}{M. Artin}
\begin{center}
M. Artin
\end{center}
\phantomsection\MakeLinkTarget*{sga4.exp.XIII}\label{sga4.exp.XIII}

\setcounter{pageoriginal}{131}

\section{射影且平坦的情形}\etiquette[1]{XIII.sec1}\pageoriginale

回顾 \hyperrefext[XII][XII.thm5.9]{5.9 bis} 在此处所考虑情形下的陈述：

\begin{enonce*}{命题 1.1}\etiquette[1.1]{XIII.prop1.1}
设 $S$ 为 Noether Hensel 环（henselian ring）的谱，
$f:X\rightarrow S$ 为射影且平坦的态射（projective and flat morphism）。设 $X_{\circ}$ 为
$X/S$ 的闭纤维。限制函子
$$
\Et(X)\longrightarrow \Et(X_{\circ})
$$
是范畴等价。
\end{enonce*}

如前所述，由 \hyperrefext[XII][XII.prop6.5]{6.5} (i) 可知该函子完全忠实。
还须证明，$X_{\circ}$ 的每个平展覆盖（étale covering）$Y_{\circ}$ 都由
$X$ 的某个平展覆盖 $Y$ 诱导。

\begin{enonce*}{引理 1.2}\etiquette[1.2]{XIII.lem1.2}
设 $S$ 局部 Noether，$f:X\rightarrow S$ 为射影且平坦的态射，
$Y\subset X$ 为子概形。终函子的子函子
$\epsilon=\prod_{X/S}Y/X$ 
“表示条件 $Y=X$”，即对 $S'\rightarrow S$ 有
$$
\epsilon(S')=
\begin{cases}
\phi & \mbox{若 }\quad Y'\neq X',\\
\{\phi\} &\mbox{若 }\quad Y'=X',
\end{cases}
$$
（其中\pageoriginale 撇号表示基变换（base change）
$S'\rightarrow S$ 的结果），可由 $S$ 的一个子概形 $Z$ 表示。若
$Y$ 是开的（\resp 闭的），则 $Z$ 亦然。
\end{enonce*}

\noindent
\emph{证明}。设 $Y$ 为 $X$ 的某个开子集 $U$ 中的闭子概形，
$C$ 为 $U$ 的闭补集。由于 $f$ 固有（proper），
$C$ 在 $S$ 中的像为闭子集 $D$，且显然条件
$U=X$ 由 $S$ 的开子集 $S-D$ 表示。因此可假设
已有 $U=X$，即 $Y$ 是闭子概形。
考虑 $\Oo_X$-模的满射态射
$\Oo_X\xrightarrow{\alpha}\Oo_Y$；条件 $Y=X$
等价于存在态射
$\Oo_Y\xrightarrow{\beta}\Oo_X$，使复合 $\beta\alpha$ 为
恒等态射，即存在 $\Hhom(\Oo_Y,\Oo_X)$ 的一个截面，
提升 $\Hhom(\Oo_X,\Oo_X)$ 的恒等截面。
由 \EGA III 7.7.8 可知，将 $S'/S$ 对应到
$\Hhom(\Oo_Y,\Oo_X)$（\resp $\Hhom(\Oo_X,\Oo_X)$）的
截面集的函子，可由 $S$ 上的向量丛（vector bundle）$V_1$（\resp $V_2$）表示。
设 $\varphi:V_1\rightarrow V_2$ 为由 $\Oo_X\rightarrow\Oo_Y$ 诱导的态射
$u\mapsto u\alpha$，$s:S\rightarrow V_2$ 为
对应于恒等态射 $\Oo_X\rightarrow\Oo_X$ 的截面。
截面 $s$ 在 $\varphi$ 下的逆像可与 $S$ 的一个闭子概形
$Z$ 等同，且显然 $Z$ 表示函子 $\epsilon$，
引理得证。

\begin{enonce*}{引理 1.3}\etiquette[1.3]{XIII.lem1.3}
设 $S$ 为局部 Noether 概形，$f:X\rightarrow S$ 为
射影且平坦的态射，$E$ 为 $X$ 上的局部自由层（locally free sheaf）。
设
$$
QE:(\Sch)/S\longrightarrow(\Ens)
$$
为如下函子：它将 $S'/S$ 对应到 $E'=E\otimes_{\Oo_S}\Oo_{S'}$ 的
赋有平展 $\Oo_{X'}$-代数结构的商层 $A'$ 所成的集合 $QE(S')$。
此函子可由 $S$ 上的局部有限型概形表示。
\end{enonce*}

\noindent
\emph{证明}。\pageoriginale 函子 $\sheaf{\Quot}(E)$ 可由一个局部有限型概形
（TDTE IV，第~221 讲）表示，记为 $Z/S$；其中
$\sheaf{\Quot}(E)(S')$ 是 $E'$ 的在 $S'$ 上平坦的商所成的集合。
设 $M_u$ 为 $X_Z$ 上 $E\otimes_{\Oo_S}\Oo_Z$ 的“泛”商。
$\sheaf{\Quot}(E)$ 中由在 $X'$ \emph{上}平坦的商所成的子函子，
由 $Z$ 的一个开子集 $U$ 表示。事实上，
设 $M'$ 为 $E'$ 的一个商，由某个 $S$-态射 $S'\rightarrow Z$ 给定。
由于 $M'$ 在 $S'$ 上平坦，它在 $X'$ 上平坦等价于：
对每个点 $s'\in S'$，模 $M'_{s'}$ 在纤维 $X'_{s'}$ 上平坦
（\EGA IV 5.9）。这又等价于
$M_u|_{X_z}$ 在纤维 $X_z$ 上平坦，其中 $z$ 为
$s'$ 的像。然而，由 \EGA IV 11.1.5，使
$M_u|_{X_z}$ 在 $X_z$ 上平坦的点 $z\in Z$ 所成的集合 $U$，
也就是使 $M_u$ 在 $X_z$ 的每一点处均在 $X_Z$ 上平坦的点集，
在 $Z$ 中是开的，故断言成立。

现在只须表示相对函子 $QE/U[f]$。于是将
$S$ 换为 $U$，$X$ 换为 $\fprod{X}{U}{S}$ 后，
问题就归结为以下引理：

\begin{enonce*}{引理 1.4}\etiquette[1.4]{XIII.lem1.4}
设 $S$ 为局部 Noether 概形，$f:X\rightarrow S$ 为
射影且平坦的态射，$M$ 为局部自由
$\Oo_X$-模。函子
$\Alget:(\Sch)/S\rightarrow(\Ens)$ 将 $S'/S$ 对应到
$M'$ 上的在 $\Oo_{X'}$ 上平展的代数结构所成的集合
$\Alget(S')$；此函子可由 $S$ 上的有限型概形表示。
\end{enonce*}

\noindent
\emph{证明}。为简便起见，以下说“可表示”时，
均理解为“可由一个\emph{有限型}相对概形表示”。
下文出现的所有向量丛均由\emph{凝聚}层定义，
因而都是有限型的。$M'$ 上的一个双线性合成律，
由局部自由层 $\Hhom(M'\otimes M',M')$ 的一个截面给定。
由于 $\Hhom(M'\otimes M',M')$ 在 $S$ 上平坦，
由\pageoriginale \EGA III 7.7.6 可知，“双线性合成律”函子
由 $S$ 上的一个向量丛 $Z$ 表示。还须验证，
合成律满足结合律、交换律及具有单位截面的附加条件，
可由 $Z$ 的一个子概形表示。考虑相对函子，
即可将 $S$ 换为 $Z$，因而可假设已给定
$M$ 上的一个合成律，只须证明表示上述条件的
$S$ 的子函子可由 $S$ 的子概形表示。

对结合律而言，条件是
$\Hhom(M\otimes M\otimes M,M)$ 的两个截面在基变换
$S'\rightarrow S$ 后相等。设 $Y$ 为这两个截面在对应于
$\Hhom(M\otimes M\otimes M,M)$ 的向量丛上的纤维积。
概形 $Y$ 可与 $X$ 的一个闭子概形等同，而在
$S'\rightarrow S$ 后成立的结合律表示为条件 $Y'=X'$。
因此，由 \hyperref[XIII.lem1.2]{1.2}，它定义了一个
可由 $S$ 的闭子概形表示的 $S$ 的子函子。
对交换律而言，条件是 $\Hhom(M\otimes M,M)$ 的两个截面相等，
由 \hyperref[XIII.lem1.2]{1.2}，这同样定义一个可由闭子概形表示的
$S$ 的子函子。最后，如今可假设合成律
已经满足结合律及交换律，并考虑单位截面存在的条件；
此时这样的截面必定唯一。将 $S'/S$ 对应到
$M'$ 的截面集的函子，由 \EGA III 7.7.6 可由一个向量丛表示，
于是只须表示此向量丛的对应于\emph{单位截面}的子函子
（这里利用刚才指出的唯一性）。如常转到相对情形，
即可假设已经给定 $M$ 的一个截面 $\epsilon$，
并证明表示它变为\pageoriginale 单位截面这一条件的
$S$ 的子函子可表示。然而，$\epsilon$ 为单位的条件，
正是显然的态射的复合
$M'\rightarrow M'\otimes M'\rightarrow M'$ 为恒等态射，
即 $\Hom(M,M)$ 的两个指定截面相等。
再次由 \hyperref[XIII.lem1.2]{1.2} 即得结论。

最后，假设局部自由模 $M$ 已赋有一个含幺交换代数结构。
它平展的点所成的集合在 $X$ 中是开的（\SGA I 4.5），
因此可在 $S$ 上由一个（开）子概形表示。
这就证明了 \hyperref[XIII.lem1.4]{1.4}，
从而也证明了 \hyperref[XIII.lem1.3]{1.3}。

\setcounter{subsection}{4}
\subsection{}\etiquette[1.5]{XIII.subsec1.5}
现设 $Z/S$ 为表示 \hyperref[XIII.lem1.3]{1.3} 中所考虑函子的对象，
并考虑 $Z$ 在某点 $z$ 处在 $S$ 上光滑的条件。
将 \SGA I III 3.1 译为函子 $QE$ 的语言，
即得到如下条件：

设 $B$ 为 Artin 局部环，$\Mm=\rad B$，$J$ 为
$B$ 的满足 $J\Mm=0$ 的理想，置 $S'=\Spec B$，$S''=\Spec B/J$，
并设给定态射 $S'\rightarrow S$。对于基变换
$S'\rightarrow S$ 及 $S''\rightarrow S$ 的结果，采用显然的记号。

\begin{itemize}
\item[$(\ast)$] 对 $E''$ 的每个赋有平展
$\Oo_{X''}$-代数结构、且对应于像为 $z$ 的
$S$-态射 $S''\rightarrow Z$ 的商 $A''$，都存在
$E'$ 的一个赋有平展 $\Oo_{X'}$-代数结构的商
$A'$，使其诱导出所给的商。
\end{itemize}

由于 $X'$ 和 $X''$ 的底空间相同，诱导
$A''$ 的代数 $A'$ 已在唯一同构意义下确定
（\SGA I 8.3）。\pageoriginale 因此，可将条件
$({\ast}{\ast})$ 表述如下：

\begin{itemize}
\item[$({\ast}{\ast})$] 设 $A'$ 为 $\Oo_{X'}$ 上的平展代数，
$A''$ 为 $A'$ 在 $X''$ 上诱导的代数。每个
满同态
$$
E''\xrightarrow{\varphi''}A''
$$
若对应于像为 $z$ 的态射 $S''\rightarrow Z$，则均可提升为同态
（由 Nakayama 引理，该同态必定为满同态）
$$
E'\xrightarrow{\varphi'}A'{\quoi}.
$$
\end{itemize}

由正合列
$$
0\longrightarrow J\longrightarrow B\longrightarrow B/J\longrightarrow
0
$$
得到 $X'$ 上的层的正合列：
$$
0\longrightarrow J\otimes\Hhom(E',A')\longrightarrow
\Hhom(E',A')\longrightarrow\Hhom(E'',A'')\longrightarrow
0{\quoi}.
$$
提升 $\varphi''$ 的障碍位于
$H^1(X',J\otimes\Hhom(E',A'))$ 中。设
$S_{\circ}=\Spec B/\Mm$。由于 $J\Mm=0$，有
$$
J\otimes\Hhom(E',A')\simeq
J\otimes\Hhom(E_{\circ},A_{\circ}){\quoi},
$$
且显然
$$
H^1(X',J\otimes\Hhom(E_{\circ},A_{\circ}))\simeq
J\otimes H^1(X_{\circ},\Hhom(E_{\circ},A_{\circ})){\quoi}.
$$
考虑到上同调与基域扩张可交换，得到：

\begin{enonce*}{推论 1.6}\etiquette[1.6]{XIII.cor1.6}
设 $Z/S$ 为表示 \hyperref[XIII.lem1.3]{1.3} 中函子的概形，
置 $X_Z=\fprod{X}{Z}{S}$，并设 $A_Z$ 为
$E_Z$ 的（“一般”）平展商代数。则在每个满足
$$
H^1(X_{\circ},\Hhom(E_{\circ},A_{\circ}))=0{\quoi},
$$
的点 $z_{\circ}$ 处，$Z/S$ 都是光滑的；其中\pageoriginale 下标零
表示限制到 $X_Z/Z$ 在点 $z_{\circ}$ 处的纤维
$X_{\circ}$。
\end{enonce*}

现在可以完成 \hyperref[XIII.prop1.1]{1.1} 的证明了。采用该命题的记号，
设 $Y_{\circ}/X_{\circ}$ 为一个平展覆盖，
由 $\Oo_{X_{\circ}}$ 上平展的代数层 $A_{\circ}$ 给定。
当 $n$ 充分大时，$A_{\circ}(n)$ 由其截面生成
（\EGA III 2.2.1）。于是有满射
$$
\Oo^N_{X_{\circ}}\longrightarrow A_{\circ}(n)\longrightarrow
0{\quoi},
$$
因而有满射
$$
\Oo^N_{X_{\circ}}(-n)=E_{\circ}\longrightarrow
A_{\circ}\longrightarrow 0{\quoi}.
$$
取 $n$ 充分大，使另有
\begin{equation*}
H^1(X_{\circ},\Hhom(E_{\circ},A_{\circ}))=0{\quoi},\tag{$+$}
\end{equation*}
由 \EGA III 2.2.1，这样的选取是可能的，因为
$$
\Hhom(E_{\circ},A_{\circ})
=[\Hhom(\Oo^N_{X_{\circ}},A_{\circ})](n)
=[A_{\circ}(n)]^N{\quoi}.
$$
置 $E=\Oo^N_X(-n)$，设 $Z/S$ 为表示
\hyperref[XIII.lem1.3]{1.3} 中所考虑函子的概形。
$E_{\circ}$ 的商 $A_{\circ}$ 对应于 $Z/S$ 的闭纤维中的
一个点 $z_{\circ}\in Z$，它是 $k(s)$ 上的有理点；
由 \hyperref[XIII.subsec1.5]{1.5} 及 $(+)$，
$Z/S$ 在 $z_{\circ}$ 处光滑。由于 $S$ 为\emph{Hensel} 环的谱，
$Z/S$ 有一个经过 $z_{\circ}$ 的截面；例如，
这可直接由光滑性的原始定义（\SGA II 1.1）得到。
因而，$A_{\circ}$ 由 $X$ 上平展的某个代数 $A$
（它是 $E$ 的商）诱导，这就完成了
\hyperref[XIII.prop1.1]{1.1} 的证明。


\section{相对维数至多为一的情形}\etiquette[2]{XIII.sec2}\pageoriginale

断言如下。

\begin{enonce*}{命题 2.1}\etiquette[2.1]{XIII.prop2.1}
采用 \textup{\hyperrefext[XII][XII.thm5.9]{5.9 bis}} 的记号，若
$S$ 为 Noether 严格局部概形，且 $f$ 为相对维数
$\leq1$ 的射影态射，则函子
$\Et(X)\rightarrow\Et(X_{\circ})$ 是范畴等价。

由 \textup{\hyperrefext[XII][XII.lem8.5]{8.5}}，这将完成
\textup{\hyperrefext[XII][XII.thm5.1]{5.1} (ii)} 的证明；
后者同时推广了 \hyperref[XIII.prop1.1]{1.1} 及
\hyperref[XIII.prop2.1]{2.1}。
\end{enonce*}

\begin{enonce*}{引理 2.2}\etiquette[2.2]{XIII.lem2.2}
设 $S$ 为 Noether 局部环 $A$ 的谱，$X/S$ 为
相对维数 $\leq n$ 的射影概形。则存在相对维数
$\leq n$ 的射影且平坦的 $S$-概形 $Z/S$，以及
$S$ 上的闭浸入 $X\hookrightarrow Z$。
\end{enonce*}

\noindent
\emph{证明。}设 $X\rightarrow\PP^N_S$ 为 $S$ 上的一个射影嵌入，
$J\subset A[x_{\circ},\ldots,x_N]$ 为定义 $X$ 的
齐次理想。置 $k=A/\mathfrak m_A$，$S_{\circ}=\Spec k$，
$X_{\circ}=X\otimes_Sk$。则 $X_{\circ}$ 由齐次理想
$J_{\circ}=\Iim(J\rightarrow k[x_{\circ},\ldots,x_N])$ 定义。
由于 $\dim X_{\circ}\leq n$，众所周知，存在
$J_{\circ}$ 中的元素 $f^{\circ}_1,\ldots,f^{\circ}_{N-n}$，
定义 $\PP^N_k$ 的一个维数为 $n$ 的子概形 $Z_{\circ}$。
设 $f_i$ $(1\leq i\leq N-n)$ 为 $J$ 中提升
$f^{\circ}_i$ 的元素，$Z$ 为由
$f_1,\ldots,f_{N-n}$ 定义的 $\PP^N_S$ 的子概形。
由 \EGA IV 11.3，$Z$ 在 $S$ 上平坦；由
\EGA IV 13.1.5，其相对维数为
$\dim Z_{\circ}=n$。\hyperref[XIII.lem2.2]{2.2} 得证。

\noindent
\emph{\hyperref[XIII.prop2.1]{2.1} 的证明。}设 $X/S$ 为相对维数
$\leq1$ 的射影概形，$S$ 为 Noether 严格局部概形。
可假设 $X$（从而 $X_{\circ}$）连通。设
$X\rightarrow Z$ 为 $S$ 上的一个闭浸入，其中 $Z$ 射影、
平坦，且相对维数 $\leq1$；见
\hyperref[XIII.lem2.2]{2.2}。设 $s_{\circ}$ 为
$S$ 的闭点。有相应的\pageoriginale 范畴图式
$$
\xymatrix@C=2.5cm@R=1.25cm{
\Et(Z)\ar[r]\ar[d]^{a}&\Et(X)\ar[d]^{b}\\
\Et(Z_{\circ})\ar[r]^{c}&\Et(X_{\circ}).
}
$$
要证明 $b$ 为等价。由 \hyperrefext[XII][XII.lem5.8]{5.8}，
$b$ 完全忠实，因此只须证明它本质满射。
由 \hyperref[XIII.prop1.1]{1.1}，函子 $a$ 是等价；
因而只须证明 $c$ 本质满射。略去下标
${}_{\circ}$，问题就归结为以下命题。

\begin{enonce*}{命题 2.3}\etiquette[2.3]{XIII.prop2.3}
设 $k$ 为可分闭域，$Z$ 为维数
$\leq1$ 的局部有限型 $k$-概形，$X\hookrightarrow Z$
为连通闭子概形。则 $X$ 的每个平展覆盖都由
$Z$ 的某个平展覆盖诱导。
\end{enonce*}

\noindent
\emph{证明。}可假设 $Z$ 和 $X$ 的维数均为 $1$，
因为当 $X$ 为零维时，结论由 $k$ 的假设立即得到。
此时 $Z=X\cup Y$，其中 $Y$ 是取 $Z-X$ 的闭包所得的
闭子概形。置 $V=X\cap Y$；这是一个离散的，
即局部 Artin 的概形。态射
$X\amalg Y\rightarrow Z$ 有限且满射，故为平展覆盖范畴的
严格下降态射（strict descent morphism）
（\SGA \hyperrefext[IX][IX.cor4.7]{4.7}）。显然
\begin{equation*}
\fprod{(X\amalg Y)}{(X\amalg Y)}{Z}
=X\amalg Y\amalg V^{(1)}\amalg V^{(2)}
\quad(V^{(i)}=V){\quoi},\tag{$\ast$}
\end{equation*}
其中 $\pr_i$ 将 $V^{(i)}$ 映到 $X\amalg Y$ 的分支 $X$，
而当 $j\neq i$ 时，将 $V^{(j)}$ 映到分支 $Y$。

设 $X'/X$ 为次数 $n$ 的平展覆盖，$Y''/Y$ 为
次数 $n$ 的一个覆盖（例如次数 $n$ 的平凡覆盖），
于是 $X'\amalg Y''$ 是\pageoriginale $X\amalg Y$ 的
次数 $n$ 的覆盖。只须找出此覆盖相对于
$X\amalg Y\rightarrow Z$ 的下降资料。考察
$(\ast)$ 可知，这样的下降资料恰好等价于一个同构
$$
\fprod{X'}{V}{X}=V'\xrightarrow{\sim}
V''=\fprod{Y''}{V}{Y}{\quoi}.
$$
由于 $k$ 可分闭，$V$ 各点的剩余域亦然；
所以 $V$ 的每个平展覆盖都完全分裂。
因此 $V'\simeq V''$，结论得证。

\begin{enonce*}[remark]{注 2.4}\etiquette[2.4]{XIII.rem2.4}
可以用一个直接论证，避开 \SGA IX 中的精细结果及下降理论：
只须证明，$Z$ 的平展覆盖范畴等价于三元组
$(X',Y',\theta)$ 所成的范畴，其中
$X'$（分别 $Y'$）为 $X$（分别 $Y$）的平展覆盖，
而 $\theta$ 为一个 $V$-同构
$\fprod{X'}{V}{X}\simeq\fprod{Y'}{V}{Y}$。
对任意可表为两个闭子概形 $X$ 与 $Y$ 之并的
概形 $Z$，这一结果都成立；它可直接由
$Z$ 上拟凝聚模的类似描述证明
（此时 $X=V(\mathfrak J)$，$Y=V(\mathfrak K)$，且
$\mathfrak J\cap\mathfrak K=0$）。
\end{enonce*}

\footnotetext[1]{这一结果也以稍一般的形式出现在 \EGA IV 21.9.11 及 21.9.12 中；
特别地，在 \hyperref[XIII.cor3.2]{3.2} 中，只须要求
$f$ \emph{分离}，不必要求它固有。}
\section{关于 Picard 群的一个辅助结果 ${}^{(*)}$}\etiquette[3]{XIII.sec3}

\begin{enonce*}{命题 3.1}\etiquette[3.1]{XIII.prop3.1}
设 $X$ 为 Noether 概形，$X_{\circ}$ 为 $X$ 的一个闭子概形，
满足 \textup{(a)} $\dim X_{\circ}\leq1$，以及
\textup{(b)} 对 $X$ 的每个非空连通闭子集 $Y$，
$Y_{\circ}=Y\cap X_{\circ}$ 非空且连通。
\begin{enumerate}
\renewcommand{\theenumi}{\roman{enumi}}
\renewcommand{\labelenumi}{(\theenumi)}
\item 设 $Z_{\circ}$ 为所有仅含一个点的集合
$\overline{\{x\}}\cap X_{\circ}$ 的并，其中 $x\in\Ass\Oo_X$。
对 $X_{\circ}$ 上的每个满足 $\supp(D_{\circ})$ 不与
$Z_{\circ}$ 相交的 Cartier 除子 $D_{\circ}$，
存在 $X$ 上的一个 Cartier 除子 $D$ 诱导出 $D_{\circ}$
（且当 $D_{\circ}\geq0$ 时，可取 $D\geq0$）。

\item 假设 $X_{\circ}$ 有一个包含有限集
$\Ass\Oo_{X_{\circ}}\cup Z_{\circ}$ 的仿射开子集 $U_{\circ}$
（若 $X_{\circ}$ 有丰沛可逆层，此条件自动满足）。
则 $\Pic(X)\rightarrow\Pic(X_{\circ})$ 为满射。
\end{enumerate}
\end{enonce*}

\noindent
\emph{证明。}
\begin{enumerate}
\renewcommand{\theenumi}{\roman{enumi}}
\renewcommand{\labelenumi}{(\theenumi)}
\item 由于 $\dim X_{\circ}\leq1$，$X_{\circ}$ 上的每个
Cartier 除子 $D_{\circ}$ 显然具有离散支集，
因而可写成两个有效 Cartier 除子之差，且它们的支集均包含在
$D_{\circ}$ 的支集中。因此，(i) 归结为
$D_{\circ}$ 有效的情形。对每个
$x\in\Supp D_{\circ}$，它在 $x$ 处由一个非零因子
$f_x\in\rad(\Oo_{X_{\circ},x})$ 定义。元素 $f_x$
由 $\Oo_X$ 在 $x$ 的某个开邻域 $U_x$ 上的一个截面
$g_x$ 诱导。断言：缩小 $U_x$ 后，可使截面
$g_x$ 成为 $\Oo_{U_x}$ 中的非零因子，即
$V(g_x)\cap\Ass\Oo_{U_x}=\varnothing$。

事实上，由条件 (a) 及 (b)，对每个
$z\in\Ass\Oo_X$，集合 $\overline{\{z\}}\cap X_{\circ}$
为 $X_{\circ}$ 的连通闭子集，故或者仅含一个闭点，
或者是 $X_{\circ}$ 的若干一维不可约分支的连通并。
由于 $x\notin Z_{\circ}$，在前一种情形下有
$x\notin\overline{\{z\}}$。在后一种情形下，
不可能同时有 $x\in\overline{\{z\}}$ 和 $z\in V(g_x)$，
否则 $g_x$，从而 $f_x$，将在
$\overline{\{z\}}\cap X_{\circ}$ 的某个不可约分支上恒为零，
而该分支也是 $X_{\circ}$ 的不可约分支；
这与 $f_x$ 为 $\Oo_{X_{\circ},x}$ 中的非零因子矛盾。
因此缩小 $U_x$ 后，可\pageoriginale 假设
$V(g_x)\cap\Ass\Oo_{U_x}=\varnothing$，故 $g_x$ 为非零因子。

设 $Y_x$ 为 $V(g_x)$ 在 $X$ 中的闭包。点 $x$ 在
$Y_x\cap X_{\circ}\cap U_x$ 中孤立，因而在
$Y_x\cap X_{\circ}$ 中也孤立。由条件 (b)，$Y_x$
有一个连通分支 $Y'_x$，满足 $Y'_x\cap X_{\circ}=\{x\}$。
而且 $Y'_x\subset U_x$：事实上，
$T_x=Y'_x-(Y'_x\cap U_x)$ 在 $X$ 中闭，且满足
$T_x\cap X_{\circ}=\varnothing$，于是由条件 (b) 得
$T_x=\varnothing$。设 $D'_x$ 为 $X$ 上的 Cartier 除子，
它在 $U_x$ 上由 $Y'_x$ 诱导，在 $X-Y'_x$ 上由零除子诱导。
则 $D=\sum_xD'_x$ 为 $X$ 上的有效 Cartier 除子，
在 $X_{\circ}$ 上诱导出 $D_{\circ}$。

\item 结论由 (i) 及如下事实得到：$X_{\circ}$ 上的每个
可逆模 $\LL_{\circ}$ 都同构于某个支集不与
$Z_{\circ}$ 相交的 Cartier 除子 $D_{\circ}$ 所定义的模。
后一断言意味着，$\LL_{\circ}$ 在某个包含
$\Ass\Oo_{X_{\circ}}\cup Z_{\circ}$ 的开集 $V_{\circ}$ 上
有一个截面，且该截面在此集合的每一点均不为零。
只须将 $X_{\circ}$ 换为仿射开集 $U_{\circ}$ 后证明这一点；
在此情形下，它由熟知的事实得到：半局部概形上的可逆模
与结构层同构。最后，若 $X_{\circ}$ 有丰沛可逆层，
则 $X_{\circ}$ 的每个有限子集都包含在某个仿射开子集中。
这就证明了 (ii)。
\end{enumerate}

有如下推论，它推广了 \hyperrefext[XII][XII.prop8.6]{8.6}，
并完成了 \hyperrefext[XII][XII.thm5.1]{5.1} (iii) 的证明。

\begin{enonce*}{推论 3.2}\etiquette[3.2]{XIII.cor3.2}
设 $S$ 为 Noether Hensel 环的谱，
$f:X\rightarrow S$ 为相对维数 $\leq1$ 的固有态射。
对 $X$ 的每个底空间与 $X/S$ 的闭纤维\pageoriginale
$X_{\circ}$ 相同的闭子概形 $X'_{\circ}$，典范映射
$\Pic(X)\rightarrow\Pic(X'_{\circ})$ 均为满射。
\end{enonce*}

事实上，将命题应用于 $(X,X'_{\circ})$。
条件 (a) 由假设成立，条件 (b) 由
\hyperrefext[XII][XII.cor5.7]{5.7} 得到。还须验证
(ii) 所要求的 $U_{\circ}$ 存在；这是因为
$X'_{\circ}$ 必为射影的，正如域上的每条固有曲线都是射影的一样。

\begin{enonce*}[remark]{注}
若读者不愿假定或重建域上的固有曲线均为射影这一事实，
可以只在 $f$ 为射影态射这一附加假设下使用此推论；
这对于此处所需的应用已经足够。反之，
若承认域上的固有代数曲线的射影性，则可得出：
在 \hyperref[XIII.cor3.2]{3.2} 的假设下，$X$
必在 $S$ 上射影。事实上，取一个丰沛元素
$\xi_{\circ}\in\Pic(X_{\circ})$，将它提升为某个
$\xi\in\Pic(X)$，再应用 \EGA III 4.7.1 即可。
\end{enonce*}

\begin{thebibliography}{99}
\bibitem{sga4.exp.XIII.bib.TDTE} Grothendieck A., Technique de descente et
théorèmes d'existence en Géométrie Algébrique, in Fondements de la
Géométrie Algébrique (Extraits du Sém. Bourbaki 1957--62), Secrétariat
mathématique, 11, rue Pierre Curie.
\end{thebibliography}

